\chapter{数学与博弈论工具}

数字经济的理论模型无论多复杂，落到推导层面都依赖三类基本工具：最优化方法求解个体的最优选择，博弈论刻画多主体的策略互动，概率与统计处理不确定性与信息。本章把这套工具系统梳理一遍，为第 3 章起的模型章节提供技术支撑。

\begin{examguide}
\textbf{本章考情}：本章为工具章，内容对应考研数学三与 396 经济类联考的微积分、概率论基础，以及专业课中博弈论考题的常用工具（\important{基础}）。本章不单独命题，但全书后续各章的推导均依赖本章结论：古诺模型与反应函数（第 4 章平台竞争）、逆向归纳（第 6 章自我优待）、重复博弈（第 6 章算法合谋）、协调博弈（第 3 章标准竞争）、拍卖理论（第 7 章数据定价）、贝叶斯更新（第 5 章基于消费历史的定价）、信息甄别（第 10 章大数据风控）。

\textbf{核心考点}：（1）拉格朗日乘数法、包络定理与隐函数定理；（2）占优策略、严格纳什均衡与纳什均衡的定义及求解：划线法、反应函数法与混合策略无差异条件；（3）动态博弈与重复博弈：逆向归纳、子博弈精炼均衡、斯塔克尔伯格均衡、触发策略临界值 $\delta^*$；（4）信息经济学：柠檬市场、逆向选择的保险市场与信息甄别、委托代理、信号传递与拍卖理论；（5）期望、方差、条件概率，贝叶斯更新与期望效用。

\textbf{本章主线}：最优化方法（2.1）→ 博弈论（2.2）→ 信息经济学（2.3）→ 概率与统计（2.4）。
\end{examguide}

\section{最优化方法}

经济学对"选择"的处理遵循统一的范式：把消费者、厂商、平台等行为主体刻画为在约束下追求某个目标的\textbf{最优化问题}。最优化研究两类基本问题：其一，给定目标函数与约束条件，最优选择是什么；其二，当外部环境（价格、收入、政策参数）变化时，最优选择与最优值如何随之变化。三个工具各司其职：拉格朗日乘数法负责求解最优选择，包络定理负责度量最优值随参数的变化，隐函数定理负责度量最优选择随参数的变化（比较静态）。\mn{最优化三问：怎么求最优解（拉格朗日）？参数变了最优值变多少（包络定理）？参数变了最优选择变多少（隐函数定理）？}三者构成最优化方法的完整链条，也是全书一切推导的技术起点。

\begin{definition}[最优化问题（optimization problem）]
在约束条件界定的可行集上，寻找使目标函数取得最大（或最小）值的选择变量的过程。
\end{definition}

\subsection{拉格朗日乘数法}

约束最优化是全书模型的共同形式：消费者在预算约束下最大化效用，厂商在技术约束下最小化成本，平台在双边需求约束下最大化利润。拉格朗日乘数法（Lagrange multiplier method）的用途是求解\textbf{带约束}的最优化问题：通过引入乘数，将"约束下最大化"改写为无约束问题的一阶条件，解题时无需先解出约束再代入目标函数。

\begin{derivation}{拉格朗日乘数法}{}
考虑问题
\[
\max_{x,y} \; f(x,y) \quad \text{s.t.} \quad g(x,y) = c.
\]
构造拉格朗日函数
\[
\mathcal{L}(x,y,\lambda) = f(x,y) - \lambda\bigl(g(x,y)-c\bigr),
\]
其中 $\lambda$ 为拉格朗日乘数。一阶条件为
\[
\frac{\partial f}{\partial x} - \lambda\frac{\partial g}{\partial x} = 0, \qquad
\frac{\partial f}{\partial y} - \lambda\frac{\partial g}{\partial y} = 0, \qquad
g(x,y) = c.
\]
最优解处内点的一阶条件要求各变量的边际收益与约束边际代价之比相等：$\frac{f_x}{g_x} = \frac{f_y}{g_y} = \lambda$。以上是一阶条件。
\end{derivation}

一阶条件只是必要条件：满足它的点可能是极大值点、极小值点，也可能是鞍点。确认内点极大值还需\textbf{二阶条件}：从候选点出发，沿所有不违背约束的可行方向，目标函数的二阶变化率都必须为负；也就是说，候选点在每一个可行方向上都必须向下弯曲，这样的点才是极大值点。"可行方向"的规范说法是\textbf{约束切空间}（tangent space），即与约束曲面相切、保持约束成立的微小变动方向全体；"沿每个方向二阶变化率为负"的规范表述，是拉格朗日函数的海塞矩阵在约束切空间上\textbf{负定}（negative definite）。实际操作中判定负定的工具是加边海塞行列式。\mn{加边海塞（bordered Hessian）判定：以约束的一阶偏导 $g_x$、$g_y$ 为边，拉格朗日函数的二阶偏导 $L_{xx}$、$L_{xy}$、$L_{yy}$ 为心，构成三阶加边行列式 $|\bar H|$。两变量情形下，极大值要求 $|\bar H| > 0$，极小值要求 $|\bar H| < 0$；变量更多时依次检验各阶加边主子式，符号正负交替为极大。}本书涉及的优化问题大多满足凹性要求，二阶条件自动成立，后续推导不再逐一检验。

乘数 $\lambda$ 本身有明确的经济含义：\important{它等于约束放松一单位时目标函数的增量}，即约束的\textbf{影子价格}。消费者问题中 $\lambda$ 是收入的边际效用，成本最小化问题中 $\lambda$ 是边际成本。这一解释来自下文的包络定理，乘数由此从求解的中间产物变为可作福利解释的经济变量。
\subsection{包络定理与比较静态}

包络定理解决的是最优化问题的第二类问题：当参数变化时，\textbf{最优值}（而非最优选择）如何变化。\mn{乘数速记：$\lambda$ = 约束的影子价格。消费者问题中是收入的边际效用，厂商成本问题中是边际成本。}直接计算需要先求出最优解随参数调整的幅度（$dx^*/dc$、$dy^*/dc$），再代入目标函数求导，步骤繁琐；包络定理表明这一步可以省去：由于最优解处一阶条件成立，最优解调整对值函数的贡献为零，参数变化只通过其直接效应起作用。

\begin{derivation}{包络定理}{}
设 $V(c) = f(x^*(c), y^*(c), c)$ 为问题
\[
\max_{x,y} f(x,y,c) \quad \text{s.t.} \quad g(x,y) = c
\]
的值函数。包络定理（envelope theorem）断言：值函数对参数的导数等于拉格朗日函数对该参数的偏导，
\[
\frac{dV}{dc} = \frac{\partial \mathcal{L}}{\partial c}\Big|_{x^*,y^*}.
\]
证明只需链式法则与一阶条件，分三步。\textbf{第一步}，对值函数求全导：
\[
\frac{dV}{dc} = f_x\frac{dx^*}{dc} + f_y\frac{dy^*}{dc} + f_c.
\]
\textbf{第二步}，代入内点一阶条件 $f_x = \lambda g_x$、$f_y = \lambda g_y$，将前两项改写：
\[
\frac{dV}{dc} = \lambda\Bigl(g_x\frac{dx^*}{dc} + g_y\frac{dy^*}{dc}\Bigr) + f_c.
\]
\textbf{第三步}，利用约束恒等式。最优解始终满足 $g(x^*(c), y^*(c)) \equiv c$，两边对 $c$ 求导得
\[
g_x\frac{dx^*}{dc} + g_y\frac{dy^*}{dc} = 1,
\]
代入即得 $dV/dc = \lambda + f_c$，恰好等于 $\partial\mathcal{L}/\partial c = f_c + \lambda$，定理得证。当目标函数不直接依赖参数时（$f_c = 0$，本书大多数应用属于此类：参数只进入约束），结论简化为
\[
\frac{dV}{dc} = \lambda^*.
\]
下面用一个标准算例检验这一结论。\textbf{算例：效用最大化。}设消费者的效用函数为 $u(x,y) = xy$，预算约束为 $p_x x + p_y y = m$，其中 $p_x$、$p_y$ 为两种商品价格，$m$ 为收入。拉格朗日函数为 $\mathcal{L} = xy + \lambda(m - p_x x - p_y y)$，一阶条件为
\[
y = \lambda p_x, \qquad x = \lambda p_y, \qquad p_x x + p_y y = m.
\]
前两式相除得 $y/x = p_x/p_y$（边际替代率等于价格比），与预算约束联立解得
\[
x^* = \frac{m}{2p_x}, \qquad y^* = \frac{m}{2p_y}, \qquad \lambda^* = \frac{m}{2p_x p_y}.
\]
将最优解代回效用函数，得到\textbf{间接效用函数}（indirect utility function）$V(m) = x^* y^* = m^2/(4p_x p_y)$。收入 $m$ 只进入约束、不进入目标函数（$f_c = 0$），由包络定理，$dV/dm$ 应等于 $\lambda^*$；直接对值函数求导，
\[
\frac{dV}{dm} = \frac{m}{2p_x p_y} = \lambda^*,
\]
两法结果一致。算例同时给出了乘数的经济含义：\important{消费者问题中的乘数等于收入的边际效用}，收入增加一单位，最大效用恰好增加 $\lambda^*$，且这一增量无需先求 $dx^*/dm$、$dy^*/dm$ 再代入，体现了包络定理"跳过最优解调整、直接看参数效应"的省力之处。
\end{derivation}

包络定理的直觉在于：当参数变化时，最优解的一阶调整对值函数的贡献为零，因为最优解处变量已处于"无改进空间"的位置，调整变量的二阶效应可忽略，参数变化只通过其直接效应起作用。这一结论使得福利分析（值函数如何随政策参数变化）变得极为简洁：无需知道最优解如何变化，只需看乘数。

先澄清两个术语。\textbf{比较静态}（comparative statics）研究的问题是：当外生参数（价格、税率、佣金率等）变化时，内生变量的均衡值如何变化。"静态"指只比较变化前后的均衡，不追究从旧均衡到新均衡的调整过程。\textbf{隐函数定理}（implicit function theorem）是实现比较静态的标准工具，它回答的是：由方程隐含定义的函数是否存在、其随参数的变化率是多少。

隐函数定理的内容可表述为：若方程 $F(x,a)=0$ 在点 $(x^*,a^*)$ 附近满足 $F_x \neq 0$，则该方程在局部将 $x$ 确定为 $a$ 的函数，且
\[
\frac{dx^*}{da} = -\frac{F_a}{F_x}.
\]
公式的直觉是：均衡条件 $F=0$ 必须始终成立，当参数 $a$ 变动 $da$ 时，内生变量 $x$ 必须调整 $dx^*$ 以抵消参数变动对 $F$ 的冲击，抵消比例恰为两偏导之比（取负号）。\mn{比较静态公式 $dx^*/da = -F_a/F_x$ 是全书均衡分析的第一工具：第 4 章平台定价、第 6 章市场界定均以此法推导。}一个具体例子：供求均衡条件 $Q_d(P) = Q_s(P,t)$（$t$ 为税率）对 $t$ 求导，即得税率变动对均衡价格的传导方向。全书各章的均衡条件均可写成 $F(\cdot)=0$ 的形式，因此这一公式是对均衡作比较静态分析的统一方法：例如第 5 章平台佣金分析中，对均衡条件关于佣金率求导，即得佣金变动对均衡价格的传导。

\section{博弈论基础}

\subsection{帕累托效率}

博弈论研究多个主体的策略互动，理解其结论首先需要一个效率基准。经济学给出的基准是帕累托效率。

\begin{definition}[帕累托改进与帕累托最优（Pareto improvement / Pareto optimality）]
若一种变动使至少一个参与人的境况变好而不使任何人的境况变坏，则称其为\textbf{帕累托改进}；当不存在任何帕累托改进的余地时，该结果达到\textbf{帕累托最优}（又称帕累托有效）。
\end{definition}

这一概念的要点是"无人受损"：帕累托改进要求每一方都至少不比原来差，因而回避了"人与人之间效用如何比较"的难题，成为无需价值判断即可达成的效率共识。\mn{帕累托改进=有人变好、无人变坏；帕累托最优=再无改进余地。纳什均衡不保证帕累托有效，囚徒困境是最简例证。}博弈论的核心张力正在于此：\important{策略互动的结果常常不是帕累托最优的}。经典例子是下文的囚徒困境：双方都背叛是纳什均衡，但（合作，合作）对双方都严格更好。均衡与效率的这种背离，正是理解平台间价格战、广告军备竞赛等现象的基本线索。

\subsection{博弈的要素与表示}

博弈由三个要素构成：\textbf{参与人}（players）、\textbf{策略}（strategies）与\textbf{支付}（payoffs）。支付是各参与人在每一策略组合下获得的效用或利润。策略是一份\textbf{完整的行动计划}：规定参与人在博弈的每一个可能情形下采取什么行动。下棋时，单独一步棋只是一个动作，"无论对手怎么走，我接下来如何应对"的全套预案才是策略。

\begin{definition}[纯策略与混合策略（pure strategy / mixed strategy）]
\textbf{纯策略}是确定地选择某一个行动的策略；\textbf{混合策略}是以概率分布随机选择多个行动的策略。
\end{definition}

纯策略可视为混合策略的退化情形（以概率 1 选择某一行动）。本节的划线法与反应函数法求解的是纯策略均衡，混合策略均衡在 2.2.4 节单独讨论。博弈的表示有两种等价形式：\textbf{标准式}（策略式）以矩阵列出所有策略组合的支付，适用于同时行动博弈；\textbf{扩展式}以博弈树表示行动的先后顺序与信息结构，适用于动态博弈。两种表示在数学上等价，区别在信息含义：标准式隐含同时行动（或行动前不知对方选择），扩展式显式刻画行动顺序。

\subsection{纳什均衡}

博弈的要素与表示回答"博弈长什么样"，纳什均衡回答"博弈的结果是什么"。

\begin{definition}[纳什均衡（Nash equilibrium）]
在策略组合 $s^* = (s_1^*, \ldots, s_n^*)$ 中，若对每个参与人 $i$ 与其任意策略 $s_i$，均有
\[
u_i(s_i^*, s_{-i}^*) \geq u_i(s_i, s_{-i}^*),
\]
则称 $s^*$ 为纳什均衡。
\end{definition}

定义中的关键词是"给定他人策略"：纳什均衡要求的是\textbf{双向的}最优反应：在他人选择不变的条件下，无人愿意单方面偏离。这一概念的学理含义在于，它把"均衡"解释为一种\important{自我实现的稳定状态}：若所有人预期该策略组合发生，则没有人有动机违背预期。\mn{纳什均衡记忆要点：给定他人策略的最优反应；双向最优、无单方偏离。}应当注意，纳什均衡不保证结果是帕累托有效的，囚徒困境中相互背叛是纳什均衡，却不是合作解。这一"个体理性与集体理性冲突"的结构在平台竞争中反复出现（第 6 章算法合谋的竞争损害分析）。

求解纳什均衡时，第一步是检查是否存在\textbf{占优策略}（dominant strategy）：无论其他参与人选择什么，都严格优于自己其他策略的策略。若某参与人拥有占优策略，直接选它即可，无需猜测对方的选择；囚徒困境中"背叛"对双方都是占优策略，因此（背叛，背叛）无需划线即可判定。占优策略均衡必然是纳什均衡，反之不成立：多数博弈中参与人没有占优策略，最优选择必须与对手的策略互相匹配，这正是划线法与反应函数法的适用场景。与之相关的一个更强概念是\textbf{严格纳什均衡}（strict Nash equilibrium）：定义中的不等号对任何 $s_i \neq s_i^*$ 都严格成立。严格均衡对支付的微小扰动稳健，本书模型解出的均衡（如古诺均衡）大多属于严格均衡。

\textbf{纯策略纳什均衡的求解。}有限博弈的纯策略纳什均衡可用划线法求解，其操作分三步：\textbf{第一步}，固定参与人 2 的每个策略，找出参与人 1 的最优反应，在其支付数字下划线（逐列划线）；\textbf{第二步}，固定参与人 1 的每个策略，找出参与人 2 的最优反应，在其支付数字下划线（逐行划线）；\textbf{第三步}，检查各策略组合：\important{两个支付均被划线的格子即为纳什均衡}，该组合中双方都处于各自的最优反应上，无人愿意单方面偏离。连续策略博弈则用反应函数法求解，其标准范例是古诺模型（Cournot model）。

\begin{figure}[htbp]
\centering
\begin{tikzpicture}[>=Stealth]
% 支付矩阵外框与分隔线
\draw[thick, second] (0,0) rectangle (9.6,4.9);
\draw[second] (0,3.6) -- (9.6,3.6);
\draw[second] (0,2.8) -- (9.6,2.8);
\draw[second] (1.4,1.4) -- (9.6,1.4);
\draw[second] (1.4,0) -- (1.4,4.9);
\draw[second] (2.6,0) -- (2.6,4.9);
\draw[second] (6.1,0) -- (6.1,3.6);
% 表头：参与人 2（列玩家）及其列标签
\node[font=\small\bfseries, color=second] at (5.5,4.25) {参与人 2};
\node[font=\small, color=structurecolor] at (4.35,3.2) {合作};
\node[font=\small, color=structurecolor] at (7.85,3.2) {背叛};
% 左侧：参与人 1（行玩家，合并格竖排）及其行标签
\node[font=\small\bfseries, color=second, rotate=90] at (0.7,1.4) {参与人 1};
\node[font=\small, color=structurecolor] at (2.0,2.1) {合作};
\node[font=\small, color=structurecolor] at (2.0,0.7) {背叛};
% 支付（每格第一数字=参与人 1，第二数字=参与人 2）
% (合作,合作) = (3,3)
\node at (3.6,2.1) {$3$}; \node at (5.1,2.1) {$3$};
% (合作,背叛) = (0,5)：参与人 2 的 5 划线
\node at (7.1,2.1) {$0$}; \node at (8.6,2.1) {$\underline{5}$};
% (背叛,合作) = (5,0)：参与人 1 的 5 划线
\node at (3.6,0.7) {$\underline{5}$}; \node at (5.1,0.7) {$0$};
% (背叛,背叛) = (1,1)：双方划线 → 纳什均衡
\node[color=winered] at (7.1,0.7) {$\underline{1}$}; \node[color=winered] at (8.6,0.7) {$\underline{1}$};
% 纳什均衡标注
\draw[winered, thick, rounded corners] (6.25,0.15) rectangle (9.45,1.25);
\node[font=\small, color=winered, anchor=north] at (7.85,-0.12) {唯一纳什均衡（背叛，背叛）};
\end{tikzpicture}
\caption{划线法求解囚徒困境}
\label{fig:underline}
\end{figure}

图~\ref{fig:underline} 以囚徒困境（prisoner's dilemma）演示划线法的完整操作：参与人 2 选合作时，参与人 1 的最优反应是背叛（$5>3$）；参与人 2 选背叛时，参与人 1 的最优反应仍是背叛（$1>0$）；对称地，参与人 2 在两种情形下的最优反应也都是背叛。唯一被双方同时划线的格子是（背叛，背叛）。\important{划线法的结果同时呈现了囚徒困境的核心命题：均衡结果是帕累托劣的}（$(1,1)$ 严格劣于 $(3,3)$），个体理性与集体理性的冲突在图上直接可见。

\begin{derivation}{古诺模型}{}
两家厂商生产同质产品，反需求函数 $P = a - Q$，$Q = q_1 + q_2$，边际成本均为 $c$。厂商 $i$ 的利润为
\[
\pi_i = (a - q_1 - q_2 - c) q_i.
\]
对 $q_i$ 求导得一阶条件：
\[
a - 2q_1 - q_2 - c = 0, \qquad a - q_1 - 2q_2 - c = 0.
\]
将 $q_1$ 表示为 $q_2$ 的函数得反应函数：
\[
q_1 = \frac{a - c - q_2}{2}, \qquad q_2 = \frac{a - c - q_1}{2}.
\]
联立求解得\important{古诺均衡}：
\[
q_1^* = q_2^* = \frac{a-c}{3}, \qquad P^* = \frac{a + 2c}{3}.
\]
反应函数还可用于刻画市场的调整过程。设厂商 2 的初始产量为 $q_2^{(0)}$，两家厂商交替取最优反应：厂商 1 对 $q_2^{(0)}$ 作出最优反应
\[
q_1^{(1)} = \frac{a - c - q_2^{(0)}}{2},
\]
厂商 2 再对 $q_1^{(1)}$ 作出最优反应，如此往复。反应函数的斜率为 $-1/2$，绝对值小于 1，每一轮迭代都把产量对均衡值的偏离压缩一半，产量序列收敛于 $(q_1^*, q_2^*)$。
\end{derivation}

反应函数斜率为负，其经济含义是：对方增产，自己的最优产量下降，因为对方增产压低了价格，使自己的边际收益下降。两条反应函数的交点即均衡：在交点处，双方都处于各自的最优反应上，无人有动机单方面改变产量。交替最优反应的收敛为均衡提供了动态解释：即使市场从任意出发点开始，厂商的交替调整也会把产量带回均衡，这一调整过程与第 3 章网络外部性市场的预期调整在结构上相同。古诺均衡是第 4 章平台竞争模型的基准：平台竞争在数学结构上与古诺模型同构，只不过策略变量从产量换成了价格或用户规模。


\subsection{混合策略纳什均衡}

前两小节求解的都是纯策略均衡，但并非所有博弈都存在纯策略均衡。\mn{混合策略求解套路：令对手各纯策略的期望支付相等（无差异条件），解出概率。}经典例子是"猜硬币"博弈：甲出硬币正反面，乙猜正面或反面，猜中者赢。若甲的策略可被预判（总是出正面），乙的最优反应（猜正面）必使甲受损，反之亦然。任何确定性的选择都会被对手利用，因此不存在双方同时最优的纯策略组合。此时纳什定理保证：\important{混合策略均衡一定存在}。混合策略的经济学含义是\textbf{以不可预测性保障自身}：随机化的目的正是让对手无法预判。猜硬币博弈的混合均衡是双方各以 $1/2$ 概率出正面：任何一方偏离该概率，另一方便可相应调整猜法获利。

\begin{derivation}{监督博弈与混合策略}{}
考虑经典的监督博弈：工人选择偷懒或努力，老板选择监督或不监督。设工人偷懒、老板不监督时工人获利 $g>0$；偷懒被抓住时受罚 $-f$；监督有成本 $c$。设老板监督概率为 $p$，工人偷懒概率为 $q$。工人偷懒的期望支付为
\[
U_{\text{偷懒}} = p\cdot(-f) + (1-p)\cdot g = g - p(f+g),
\]
努力的期望支付为 $0$。均衡要求工人对偷懒与努力无差异（若偷懒严格更优则 $q=1$，若严格更劣则 $q=0$，均与混合均衡矛盾），即
\[
g - p^*(f+g) = 0 \;\Rightarrow\; p^* = \frac{g}{f+g}.
\]
同理，老板对监督与不监督无差异：
\[
q^*\cdot(f - c) + (1-q^*)\cdot(-c) = q^*\cdot 0 \;\Rightarrow\; q^* = \frac{c}{f}.
\]
\end{derivation}

混合策略均衡的逻辑是\important{无差异原理}：参与人之所以随机化，是因为只有在对手的随机化概率使其各纯策略支付相等时，自己才愿意随机化。$p^* = g/(f+g)$ 的含义直观：偷懒的诱惑 $g$ 越大、惩罚 $f$ 越小，老板的监督概率就必须越高。这一结构在数字市场监管中有直接对应：平台违规收益越高、罚则越轻，监管抽查强度就必须越大（第 6、12 章）。
\subsection{动态博弈与逆向归纳}

动态博弈中，先行动者的策略选择必须考虑后行动者的最优反应。求解工具是逆向归纳法（backward induction）：从博弈树的末端向前推演，每一步都以后续阶段的最优反应为给定。\mn{古诺与斯塔克尔伯格对比（计算题常考）：古诺产量各 $(a-c)/3$；先动者产量 $(a-c)/2$，恰为垄断产量，后动者为其一半。}博弈树（扩展式）的读法：结点表示决策时点，枝表示可选行动，末端括号内依次为各参与人的支付（先行动者在前）。

\begin{figure}[htbp]
\centering
\begin{tikzpicture}[>=Stealth,
    node/.style={circle, draw=second, thick, fill=second!10, inner sep=1.5pt, font=\small},
    payoff/.style={font=\small},
    branch/.style={font=\small, color=structurecolor, midway, fill=white, inner sep=1pt}]
% 根结点：进入者
\node[node] (root) at (0,0) {};
\node[font=\small, anchor=east] at (-0.25,0) {进入者};
% 分支
\draw[thick, ->] (root) -- node[branch, above] {进入} (2.6,1.1);
\draw[thick, ->] (root) -- node[branch, below] {不进入} (2.6,-1.1);
% 在位者决策结点
\node[node] (inc) at (2.6,1.1) {};
\node[font=\small, anchor=south east] at (2.35,1.3) {在位者};
\draw[thick, ->] (inc) -- node[branch, above] {默许} (5.2,1.9);
\draw[thick, ->] (inc) -- node[branch, below] {斗争} (5.2,0.3);
% 支付
\node[payoff] at (5.6,2.2) {$(40,\;50)$};
\node[payoff] at (5.6,0.0) {$(-10,\;0)$};
\node[payoff] at (5.6,-1.1) {$(0,\;100)$};
% 逆向归纳：在位者选择默许（50>0），加粗该枝
\draw[main, ultra thick] (inc) -- (5.2,1.9);
\node[font=\small, color=main, anchor=north] at (3.9,1.75) {逆向归纳：选默许};
\end{tikzpicture}
\caption{进入威慑博弈}
\label{fig:gametree}
\end{figure}

图~\ref{fig:gametree} 以进入威慑博弈演示逆向归纳。潜在进入者考虑是否进入在位者的市场；若进入，在位者选择默许（分利，支付 $50$）或斗争（价格战，双方受损）。逆向归纳从末端开始：在位者在"进入"之后的最优反应是默许（$50 > 0$）；预期到这一点，进入者比较"进入得 $40$"与"不进入得 $0$"，选择进入。博弈树的求解顺序与博弈的行动顺序\textbf{相反}，逆向归纳因此得名。值得注意：若在位者能够\textbf{可信地}威胁斗争（如预先投资过剩产能使斗争成本下降），博弈结果将改写，"承诺"在动态博弈中的关键作用（第 3 章锁定、第 4 章排他性协议的分析基础）由此可见。

逆向归纳求出的均衡有一个规范名称：\textbf{子博弈精炼纳什均衡}（subgame perfect equilibrium）。\textbf{子博弈}（subgame）指从某个决策结起直至末端的部分博弈；若一个策略组合在原博弈与每个子博弈上都构成纳什均衡，则称其为子博弈精炼的。这一定义的作用是剔除\important{不可信的威胁}：上例中"进入后斗争"虽然在位者事前可以宣称，但在"进入之后"这个子博弈中不构成最优反应，因而含有斗争威胁的策略组合不是子博弈精炼均衡。纳什均衡只要求策略在事前最优，子博弈精炼额外要求每一时点的计划都经得起事后检验，逆向归纳正是求子博弈精炼均衡的标准方法。

\begin{derivation}{斯塔克尔伯格模型（Stackelberg model）}{}
在古诺模型的设定下，令厂商 1 先选择产量，厂商 2 观察到 $q_1$ 后再选择 $q_2$。逆向归纳：先求后行动者厂商 2 的最优反应
\[
q_2(q_1) = \frac{a - c - q_1}{2},
\]
再代入先行动者厂商 1 的利润最大化：
\[
\pi_1 = \left(a - q_1 - \frac{a-c-q_1}{2} - c\right) q_1 = \frac{(a-c-q_1)q_1}{2}.
\]
一阶条件给出
\[
q_1^* = \frac{a-c}{2}, \qquad q_2^* = \frac{a-c}{4}, \qquad Q^* = \frac{3(a-c)}{4}.
\]
与古诺均衡相比：$q_1^* > q_{\text{古诺}} > q_2^*$，总产量更高、价格更低。
\end{derivation}

均衡结构显示\important{先行动者获得优势}，其来源是\textbf{承诺}（commitment）：厂商 1 的产量一旦确定便构成可信承诺，厂商 2 只能在其剩余市场空间中最大化。这一结构是理解平台先动战略（第 3 章临界规模的抢跑逻辑、第 4 章排他性协议）的方法基础。逆向归纳的适用前提是各阶段的信息完备与理性共识，若先行动者无法承诺（产量可事后调整），则结果退回古诺均衡。

\subsection{重复博弈}

一次性博弈中的不合作结果，在重复博弈（repeated game）中可能被合作结果取代。\mn{$\delta^*$ 的记忆结构：分子 $T-R$ 是"背叛的诱惑"，分母 $T-P$ 是"背叛的总代价"。}设囚徒困境阶段博弈的支付为：合作得 $R$，背叛得 $T$（$T>R$），相互背叛得 $P$，被背叛得 $S$。

\begin{derivation}{重复博弈与触发策略}{}
考虑无限次重复博弈，双方采用触发策略（trigger strategy）：第一期合作，此后若对方曾背叛则永远背叛。若一方在某一期背叛，当期得 $T$，此后每期得 $P$；若坚持合作，每期得 $R$。设贴现因子（discount factor）为 $\delta \in (0,1)$，合作的价值为
\[
V_{\text{合作}} = R + \delta R + \delta^2 R + \cdots = \frac{R}{1-\delta}.
\]
背叛的价值为
\[
V_{\text{背叛}} = T + \delta P + \delta^2 P + \cdots = T + \frac{\delta P}{1-\delta}.
\]
合作可持续的条件为 $V_{\text{合作}} \geq V_{\text{背叛}}$，即
\[
\frac{R}{1-\delta} \geq T + \frac{\delta P}{1-\delta}
\;\Longleftrightarrow\;
\delta \geq \frac{T - R}{T - P} \equiv \delta^*.
\]
\end{derivation}

合作条件的含义是：合作能否维持取决于耐心程度 $\delta$ 是否超过临界值 $\delta^*$。$\delta$ 度量未来相对于现在的重要性，市场预期稳定、利率低时 $\delta$ 高，背叛的短期诱惑被未来的合作收益压倒，合作自发维持。数字经济中的直接应用是\textbf{算法合谋}（第 6 章）：定价算法以极高频率观察对手价格并即时惩罚降价，"触发"几乎是即时的，对应 $\delta \to 1$，因此即便没有任何明示协议，算法间也可能收敛到垄断价格。这是数字市场反垄断面临的全新挑战。

与囚徒困境不同，网络产业中更常见的是\textbf{协调博弈}（coordination game）：参与者与对方采取相同行动时收益最高，单独偏离只会受损。协调博弈的均衡有多个，哪一个被选中由预期与历史决定，这是第 3 章网络启动问题与标准竞争的博弈论基础。

\begin{derivation}{协调博弈与多重均衡}{}
两个参与者同时选择行动 $A$ 或 $B$。双方选择相同时各得收益 $u > 0$，选择不同时各得 $0$。支付矩阵为

\begin{center}
\begin{tabular}{c|cc}
 & A & B \\
\hline
A & $(u,u)$ & $(0,0)$ \\
B & $(0,0)$ & $(u,u)$ \\
\end{tabular}
\end{center}

逐格检验可知，该博弈有两个纯策略纳什均衡：$(A,A)$ 与 $(B,B)$；此外存在一个混合策略均衡，双方各以 $1/2$ 的概率选 $A$，协调失败以正概率出现。引入沉没成本后，协调问题出现不对称。设参与者 1 已在行动 $A$ 上投入沉没成本 $c$，若改选 $B$ 需再付 $c$；参与者 2 尚未投入，无论选择 $A$ 还是 $B$ 都需支付采用成本 $c$。此时支付矩阵为

\begin{center}
\begin{tabular}{c|cc}
 & $A$ & $B$ \\
\hline
$A$ & $(u-c, u-c)$ & $(-c, -c)$ \\
$B$ & $(-2c, -c)$ & $(u-2c, u-c)$ \\
\end{tabular}
\end{center}

逐列检验参与者 2 的选择：参与者 1 选 $A$ 时，参与者 2 跟随 $A$ 得 $u-c$、选 $B$ 得 $-c$，跟随更优（$u > 0$ 恒成立）；参与者 1 选 $B$ 时，参与者 2 选 $A$ 得 $-c$、跟随 $B$ 得 $u-c$，同样跟随更优。给定参与者 2 必然跟随，参与者 1 比较坚守 $A$（得 $u-c$）与改选 $B$（得 $u-2c$），坚守严格更优，均衡为 $(A,A)$。注意 $A$ 在此并非参与者 1 的占优策略：占优要求对对手的所有行动都更优，而这里对 $A$ 的选择建立在对参与者 2 跟随行为的确定性预期之上。
\end{derivation}

协调博弈的难点与囚徒困境不同：囚徒困境难在抵抗背叛的诱惑，协调博弈难在猜测对方的选择。会议软件是典型场景：同事之间使用腾讯会议还是钉钉，全员使用同一款软件时效用最高，选哪一款取决于对彼此选择的预期。支付工具的推广呈现同一结构的另一面：消费者与商家决定是否转向条码支付，双方都需先付出一次性成本（下载注册、购置机具），只有双方同时转向才能实现收益，单方先动反而吃亏，因此"双方都用现金"的均衡可能长期持续；主导平台的补贴与推广的作用，正在于帮助双方同时跨过启动成本。\important{沉没成本进一步使协调产生不对称的焦点：已投入的一方成为事实上的标准制定者，后来者只能跟随，网络市场的路径依赖由此而生}。若双方都已在不同技术上各有沉没投入且 $u < c$，改换行动得不偿失，"各自为政"的协调失败成为均衡，此时需要第三方焦点（政府标准、补贴、主导厂商的兼容承诺）打破僵局。第 3 章的标准竞争（NFC 双标准）与网络启动问题可直接用这一博弈分析（第 3 章详述）。


\section{信息经济学基础}

\subsection{信息不对称的类型学}

信息经济学区分两类信息不对称（information asymmetry）：\textbf{隐藏信息}（逆向选择，adverse selection）指交易前一方拥有另一方不知道的私人信息，如卖方知道商品质量而买方不知；\textbf{隐藏行动}（道德风险，moral hazard）指交易后一方的行动不可观察，如投保后投保人的努力程度。两者的共同后果是市场无法实现完全信息下的有效配置，区别在于前者发生在签约前、作用于"选择"，后者发生在签约后、作用于"行动"。数字经济中两类问题并存：平台上的卖家质量信息不对称（逆向选择，催生评价体系）、众包与零工劳动的努力不可观察（道德风险，催生算法考核）。

\subsection{柠檬市场（market for lemons）：逆向选择的经典模型}

逆向选择的最经典模型是阿克洛夫（Akerlof，1970）的柠檬市场：当买者无法辨别商品质量时，市场将如何运行？\mn{名词解释"柠檬市场"：次品市场 + 逆向选择。结论记忆：信息不对称下市场可能完全崩溃（$p=0$）。}模型的核心结论是，信息不对称下市场可能一步步萎缩，直至完全崩溃。

\begin{derivation}{柠檬市场模型}{}
二手车市场中，卖者知道车的质量 $\theta \in [0,1]$（均匀分布），买者仅知道分布。设买者对质量为 $\theta$ 的车出价 $1.5\theta$，卖者保留价为 $\theta$。若信息对称，质量为 $\theta$ 的车以 $1.5\theta$ 成交，市场出清且有效率。若信息不对称，买者只能按期望质量出价。给定市场均衡价格 $p$，愿意出售的车满足 $\theta \leq p$，市场上车的平均质量为
\[
\bar\theta(p) = \mathbb{E}[\theta \mid \theta \leq p] = \frac{p}{2}.
\]
买者按期望质量出价 $1.5\,\bar\theta = 0.75p$。均衡要求买者出价等于市场价格：$p = 0.75p$，唯一解为 $p = 0$。
\end{derivation}

均衡的推导揭示了逆向选择的自我强化逻辑：价格下降使高质量车退出（其保留价高），市场平均质量进一步下降，买者出价更低，更多好车退出，即"劣币驱逐良币"，直至\important{市场崩溃}。数字平台对这一问题有双重作用：评价系统、担保交易、平台认证等机制部分恢复了信息对称，体现平台的信誉中介功能；平台自身也可能制造新的逆向选择空间（如虚假评价）。这一模型是理解第 4 章平台中介价值与第 6 章平台治理的理论起点。
\subsection{逆向选择的保险市场与信息甄别}

柠檬市场展示了不知情的买方只能被动接受逆向选择的后果，保险市场提供了另一幅图景：保险公司作为不知情方，可以主动设计一组合同供投保人挑选，让不同风险类型的人各自"对号入座"。这种由不知情方设计菜单、让知情方自我暴露的机制称为\textbf{信息甄别}（screening），它与下文的信号传递解决同一问题，方向恰好相反：甄别的菜单由不知情方设计，信号由知情方主动发送。\mn{甄别与信号一句话：不知情方摆菜单是甄别，知情方亮资质是信号；两者都靠类型间的成本差异实现分离。}本小节以保险市场的 Rothschild--Stiglitz 模型说明这一机制的原理与代价。

\begin{derivation}{保险市场的逆向选择与甄别}{}
考虑保险市场上的两类投保人：高风险者出险概率 $\pi_H = 0.8$，低风险者 $\pi_L = 0.2$，初始财富均为 $100$，一旦出险损失全部财富，两类人的效用函数均为 $u(W) = \sqrt{W}$。保险公司竞争性经营，合同按精算公平定价（零利润）：足额保险（保额 $100$）对高风险者的公平保费为 $80$，对低风险者为 $20$。

\textbf{第一步：混同合同不可行。}若市场只提供一份按平均风险定价的合同，保费必然落在 $20$ 与 $80$ 之间；低风险者面对的保费高于其公平价格，投保意愿下降，投保人中高风险占比上升，平均赔付率随之上升，保费只能继续上调。逆向选择的机制与柠檬市场完全同构。

\textbf{第二步：甄别菜单。}保险公司转而同时提供两份合同：$H$ 合同为足额保险，保费 $80$；$L$ 合同为部分保险，保额 $x$，保费为 $0.2x$。分离的条件是\textbf{自选择约束}（self-selection constraint）：每类投保人都不愿冒充另一类。取 $x = 10$，$L$ 合同保费为 $2$，逐项检验两类人的选择。

高风险者的三个选择：买 $H$ 合同，确定财富 $100 - 80 = 20$，效用 $\sqrt{20} \approx 4.47$；冒充买 $L$ 合同，效用为
\[
0.2\sqrt{100 - 2} + 0.8\sqrt{10 - 2} \approx 0.2 \times 9.90 + 0.8 \times 2.83 \approx 4.24;
\]
不投保，效用为 $0.2\sqrt{100} = 2$。高风险者选 $H$ 合同（$4.47 > 4.24 > 2$）。

低风险者的选择：买 $L$ 合同，效用为
\[
0.8\sqrt{100 - 2} + 0.2\sqrt{10 - 2} \approx 0.8 \times 9.90 + 0.2 \times 2.83 \approx 8.49;
\]
不投保，效用为 $0.8\sqrt{100} = 8$。低风险者选 $L$ 合同（$8.49 > 8$），更不会去买保费 $80$ 的 $H$ 合同（届时确定财富 $20$，效用仅 $4.47$）。两类人各取所需，\important{分离均衡}成立。
\end{derivation}

数值算例中最值得注意的是 $L$ 合同的保障水平：低风险者只买到 $10$ 的保额，远低于其想要的足额保险。原因在于自选择约束：保额越高，$L$ 合同对高风险者的吸引力越大，保险公司只能把低风险合同的保障压到足够低，使高风险者冒充无利可图。\important{信息不对称的代价由低风险者承担，他们面临"保险不足"}。Rothschild--Stiglitz 模型还证明，当市场上低风险者占比足够高时，任何分离菜单都可能被一份按平均风险定价的混同合同侵入，此时纯策略均衡可能不存在；这一"不存在"结果本身就是逆向选择市场脆弱性的刻画。信息甄别在数字经济中应用广泛：平台的会员分级、差异化服务套餐、起送费与配送费的组合，本质上都是让不同类型用户自选择的菜单设计；大数据风控（第 10 章）与个性化定价（第 5 章）可以视为信息甄别的技术升级。

\subsection{道德风险与委托代理}

道德风险的标准建模框架是委托代理理论：委托人设计一份契约，使代理人在参与与激励两条约束下为委托人的利益行动。这一框架是第 9 章平台劳动契约与第 12 章数据治理设计的分析基础。

\begin{derivation}{委托代理基本模型}{}
委托人（principal，风险中性）雇佣代理人（agent）完成一项任务。代理人选择努力 $e \in \{0,1\}$，产出 $y \in \{y_H, y_L\}$，高努力使高产出的概率上升：$\Pr(y_H \mid e=1) = p_1 > p_0 = \Pr(y_H \mid e=0)$。代理人付出努力的成本为 $c$，保留效用为 $\bar u$。委托人的问题是在激励相容约束（IC，incentive compatibility）与参与约束（IR，participation constraint）下设计工资契约 $w_H, w_L$：
\[
\max_{w_H,w_L} \; p_1(y_H - w_H) + (1-p_1)(y_L - w_L)
\]
s.t.
\[
p_1 w_H + (1-p_1) w_L - c \geq p_0 w_H + (1-p_0) w_L \quad \text{(IC)}
\]
\[
p_1 w_H + (1-p_1) w_L - c \geq \bar u \quad \text{(IR)}
\]
最优契约将 \important{IR 取等号}，并从 IC 解出最优工资结构。若代理人为风险中性，最优解为 $w_H = y_H - F$、$w_L = y_L - F$，其中 $F$ 为固定费用，代理人"承包"任务，成为剩余索取者；若代理人为风险厌恶，最优契约满足 $w_H > w_L$，且工资差距小于风险中性情形。
\end{derivation}

委托代理理论的中心命题是\textbf{激励与保险的权衡}：把支付与产出挂钩（激励）会把产出波动风险转嫁给代理人（风险成本），最优契约在两者间折中。数字平台劳动中的计件工资（骑手按单计价）正是强激励契约：平台以几乎完全的承包制规避道德风险，同时把需求波动的风险完全转嫁给劳动者。这一制度设计在效率与公平之间的权衡，是第 9 章平台劳动争论的理论焦点。
\subsection{信号传递（signaling）}

逆向选择可由知情方主动发送信号而缓解。\mn{信号有效的两个条件：成本与类型负相关 + 信号不可伪造。刷单、买量破坏的正是第二个条件。}信号有效的条件有二：其一，信号成本与私人信息负相关（好类型发信号的成本更低）；其二，信号必须可观察且不可轻易伪造。

\begin{derivation}{教育信号模型}{}
两类工人：高能力者（$H$，占比 $\lambda$）生产率 $y_H$，低能力者（$L$）生产率 $y_L$，$y_H > y_L$。雇主不能观察能力，只能观察教育年限 $e$。设高能力者获得教育 $e$ 的成本为 $c_H e$，低能力者为 $c_L e$，且 $c_H < c_L$。在分离均衡中，高能力者选择教育 $e^*$，低能力者选择 $0$，雇主工资为 $w(e^*) = y_H$、$w(0) = y_L$。均衡条件为两类工人各自不愿模仿对方：
\[
y_H - c_H e^* \geq y_L \quad\text{(高能力者愿意上学)}, \qquad
y_L \geq y_H - c_L e^* \quad\text{(低能力者不愿上学)}.
\]
联立得
\[
\frac{y_H - y_L}{c_L} \leq e^* \leq \frac{y_H - y_L}{c_H}.
\]
\end{derivation}

分离均衡存在的关键在于 \important{$c_H < c_L$}：同样的教育年限，高能力者承受的成本更低，教育因此可以成为区分信号，高能力者愿意"购买"足以让低能力者望而却步的信号量。教育年限 $e^*$ 本身在纯信号模型中不提高生产率，社会为信号付出了净成本。数字经济中，平台上的认证、评分等级、开源社区的贡献记录均具有信号功能；数字信号的可伪造性（刷单、买量）恰恰破坏了信号有效的第二个条件，这是第 4 章平台治理机制设计的核心问题之一。
\subsection{拍卖理论入门（auction theory）}

拍卖是数据要素定价（第 7 章）与广告位定价（第 5 章）的基础机制。\mn{四种拍卖记忆：英式=升价喊价；荷兰式=降价拍卖；一级密封=最高报价者按其报价支付；二级密封（维克瑞）=最高报价者按次高报价支付。}四种基本拍卖形式：\textbf{英式拍卖}（English auction，升价拍卖，最后出价者以其出价得标）、\textbf{荷兰式拍卖}（Dutch auction，降价拍卖，第一个接受者以其接受价得标）、\textbf{一级价格密封拍卖}（first-price sealed-bid auction，各投标人密封报价，最高报价者以其报价支付）、\textbf{二级价格密封拍卖}（second-price sealed-bid auction，又称维克瑞拍卖 Vickrey auction，最高报价者以次高报价支付）。

\begin{derivation}{最优出价策略}{}
设投标人 $i$ 对拍品的估价为 $v_i$，各估价独立同分布，分布函数 $F(\cdot)$，密度 $f(\cdot)$。\textbf{维克瑞拍卖中，诚实出价是最优策略}：设其他人的最高出价为 $b_{-i}$。若 $i$ 出价 $b_i$，其期望支付为（胜出时）$\mathbb{E}[b_{-i} \mid b_{-i} < b_i] \cdot \Pr(b_{-i} < b_i)$，收益为 $\mathbb{E}[v_i - b_{-i} \mid \text{胜出}]$。由于出价 $b_i$ 只影响胜出概率而不影响支付额（支付额由次高出价决定），任何偏离 $b_i = v_i$ 的出价只会引入两类损失：$b_i > v_i$ 时可能以超过 $v_i$ 的次高价胜出而亏损，$b_i < v_i$ 时可能输给低于 $v_i$ 的次高价而放弃正收益。故 $b_i = v_i$ 严格占优。

\textbf{一级价格密封拍卖}中，最优出价低于估价。设对称均衡出价函数为 $b(v)$（严格递增），投标人对 $b$ 求期望收益最大化：
\[
\max_b \; (v - b) \cdot \Pr(\text{胜出}) = (v - b)\, F^{n-1}\bigl(b^{-1}(b)\bigr).
\]
一阶条件（结合均衡要求 $b^{-1}(b) = v$）给出微分方程
\[
b'(v)\bigl[F(v)\bigr]^{n-1} = (v - b)(n-1)\bigl[F(v)\bigr]^{n-2} f(v),
\]
其解（均匀分布 $F(v)=v$ 于 $[0,1]$ 时）为 $b(v) = \frac{n-1}{n} v$。\important{收益等价定理}断言：在对称独立私有价值（SIPV）假设下，四种拍卖形式给卖方带来的期望收益相同。
\end{derivation}

一级价格拍卖中出价低于估价，是在"胜出概率"与"胜出盈余"之间权衡的结果：出价越高胜率越高但盈余越薄，最优出价使两者边际相等。收益等价定理的直觉是：任何拍卖机制都是对"如何将拍品配置给估价最高者、如何从投标人处提取支付"的设计，竞争程度相同的机制在期望意义上提取相同的租金。这一定理的含义深远：机制设计者更应关注机制对共谋、信息不透明等偏离假设情形的稳健性，这正是数据拍卖设计（第 7 章）中的真实难题。
\section{概率与统计基础}

\subsection{贝叶斯更新（Bayesian updating）}

概率与统计处理不确定性：当结果取决于随机事件时，决策者需要度量各种结果的可能性，并据此作出最优选择。先固定三组基本概念。\textbf{期望}（expectation）是随机变量取值的概率加权平均，离散情形为 $\mathbb{E}[X] = \sum_i x_i \Pr(X = x_i)$，连续情形以积分代替求和；\textbf{方差}（variance）$\operatorname{Var}(X) = \mathbb{E}\bigl[(X - \mathbb{E}X)^2\bigr]$ 度量取值围绕期望的离散程度，方差越大，结果的不确定性越强；\textbf{条件概率}（conditional probability）$\Pr(A \mid B) = \Pr(A \cap B)/\Pr(B)$ 刻画已知事件 $B$ 发生时事件 $A$ 的概率。

贝叶斯更新是处理不确定性的核心工具：决策者对新信息的处理，就是按贝叶斯公式用信号修正先验信念。这一工具在信号传递、平台评价体系与基于消费历史的定价（第 5 章两期模型的后验分布）中反复使用。

\begin{derivation}{贝叶斯更新}{}
设事件 $A$ 的先验概率为 $\Pr(A)$，观察到信号 $B$ 后，后验概率由贝叶斯公式（Bayes' rule）给出：
\[
\Pr(A \mid B) = \frac{\Pr(B \mid A)\,\Pr(A)}{\Pr(B)}, \qquad
\Pr(B) = \Pr(B \mid A)\Pr(A) + \Pr(B \mid \bar A)\Pr(\bar A).
\]
算例：某平台卖家分两类，优质卖家占 $20\%$。优质卖家获得好评的概率为 $95\%$，劣质卖家为 $50\%$。观察到一个好评后，该卖家为优质卖家的后验概率为
\[
\Pr(\text{优质} \mid \text{好评}) = \frac{0.95 \times 0.2}{0.95 \times 0.2 + 0.5 \times 0.8} = \frac{0.19}{0.59} \approx 0.32.
\]
\end{derivation}

贝叶斯更新刻画了理性主体如何利用新信息修正信念。上述算例的含义值得注意：一个好评仅将"优质"概率从 $20\%$ 提升到 $32\%$，远未到令人信服的水平，因为好评信号本身噪音大（劣质卖家也有一半概率得好评）。这就是平台评价体系需要海量评价才能发挥信号功能的概率论根源，也是"刷好评"有利可图的原因：少量伪造好评即可显著移动买家的后验信念。

\subsection{期望效用与风险态度}

期望效用理论是分析风险下选择的标准框架：决策者最大化各状态效用的概率加权和，而非期望收益本身。\mn{结论速记：精算公平保费（$\rho=\pi$）下，风险厌恶者全额投保 $x^*=L$。}作为其典型应用，本节推导保险需求模型：面对可投保的损失风险，消费者最优的投保额是多少。

\begin{derivation}{保险需求模型}{}
设消费者初始财富 $W$，面临概率为 $\pi$ 的损失 $L$。保险合同：保费率 $\rho$，投保额 $x$（即保费 $\rho x$，出险时赔付 $x$）。消费者的期望效用（expected utility）最大化问题为
\[
\max_x \; (1-\pi)\, u(W - \rho x) + \pi\, u\bigl(W - L - \rho x + x\bigr).
\]
一阶条件：
\[
-(1-\pi)\rho\, u'\bigl(W - \rho x^*\bigr) + \pi(1-\rho)\, u'\bigl(W - L + (1-\rho) x^*\bigr) = 0.
\]
若保险精算公平（$\rho = \pi$），则一阶条件化简为
\[
u'\bigl(W - \pi x^*\bigr) = u'\bigl(W - L + (1-\pi) x^*\bigr),
\]
由 $u'$ 严格递减得 $x^* = L$，即风险厌恶者购买全额保险。
\end{derivation}

\textbf{风险厌恶}（risk aversion，$u'' < 0$，即边际效用递减）意味着确定性的财富组合优于具有相同期望值的风险组合。精算公平的保险将不确定损失转化为确定的保费支出，\important{风险厌恶者因此愿意全额投保}。绝对风险厌恶系数 $-u''/u'$ 度量风险厌恶强度：系数越大，风险溢价越高。数字经济的应用：平台对骑手的事故保险、数据交易中的风险定价，以及隐私经济学中消费者对"隐私风险"的效用折损（第 12 章），都建立在这一框架之上。
\begin{chapterbox}
\textbf{本章考点清单}

\begin{enumerate}
\item 拉格朗日乘数法的构造与乘数的影子价格含义 \important{【基础】}
\item 包络定理的结论与证明思路 \important{【基础】}
\item 隐函数定理与比较静态：$dx^*/da = -F_a/F_x$ \important{【基础】}
\item 帕累托改进与帕累托最优，及其与纳什均衡的关系 \important{【基础】}
\item 占优策略与严格纳什均衡；纳什均衡的定义与求解（划线法、反应函数法），古诺均衡 $q^* = (a-c)/3$ 与反应函数迭代的收敛 \important{【基础】}
\item 纯策略与混合策略的区别、混合均衡的无差异条件求解（监督博弈 $p^* = g/(f+g)$）\important{【基础】}
\item 逆向归纳法、子博弈精炼均衡（剔除不可信威胁）与斯塔克尔伯格均衡（先发优势的承诺逻辑）\important{【基础】}
\item 触发策略与合作条件 $\delta \geq \delta^* = (T-R)/(T-P)$ \important{【基础】}
\item 协调博弈：双均衡结构、沉没成本焦点与支付工具启动案例（条码支付）（为第 3 章标准竞争铺垫）\important{【基础】}
\item Akerlof 柠檬市场模型：均衡崩溃的推导 \important{【基础】}
\item 逆向选择的保险市场：混同合同不可行的逻辑与 Rothschild--Stiglitz 甄别菜单（自选择约束、低风险者保险不足）\important{【基础】}
\item 委托代理模型的 IC/IR 约束与激励-保险权衡 \important{【基础】}
\item Spence 信号模型的分离均衡条件；信息甄别与信号传递的方向区别 \important{【基础】}
\item 四种拍卖形式、维克瑞拍卖的诚实出价定理、收益等价定理 \important{【基础】}
\item 期望、方差与条件概率的定义 \important{【基础】}
\item 贝叶斯更新公式与算例 \important{【基础】}
\item 期望效用、风险厌恶与保险需求 $x^* = L$ \important{【基础】}
\end{enumerate}
\end{chapterbox}

本章的工具将在后续各章中被反复调用：第 3 章的网络外部性模型是最优化方法的应用；第 4 章的平台竞争是古诺模型的多边扩展；第 6 章的算法合谋是重复博弈的延伸；第 7 章的数据拍卖直接建立在拍卖理论之上。
